\documentclass{handout}
\usepackage{handoutSetup}\usepackage{handoutShortcuts}  
\opt{bigWide}{\wideMargins}

\begin{document}
\handoutHeader

%\wideMargins\large

\begin{verbatimwrite}{\jobname.title}
Durables
\end{verbatimwrite}
%\input \jobname.title

\handoutNameMake

%\section{The Problem}

%Main question on midterm03 was on a superior model of durables spending
%volatility.  Consider using that for a future revision of this handout.


Consider a consumer who gets utility from a flow of consumption of 
nondurable goods, $c_{t}$, as well as from a stock of durable goods, 
$d_{t}$.\footnote{The basic ideas in this handout
are derived from \cite{mankiw:durgoods}.  See \cite{cdSs} for further discussion 
of the frictionless model and empirical estimates, as well as a model that incorporates transactions costs.}  The consumer's goal is to \opt{MarginNotes}{\marginpar{\tiny Define `durable.'}}
\begin{equation}
        \max \sum_{s=t}^{T} \beta^{s-t} \uFunc(c_{s},d_{s})
\end{equation}
where $d_{s}$ is the stock of the durable good, and all other 
variables are as usually defined.

We will assume that the stock of the durable good evolves over 
time according to \opt{MarginNotes}{\marginpar{\tiny Discuss alternative: `one-hoss-shay' depreciation: small prob in each period that durable completely dies.  Harder to analyze, though there is a lit that does it.}}
\begin{equation}\begin{gathered}\begin{aligned}
        d_{t+1} & =  (1-\delta)d_{t}+{x}_{t+1}, \label{eq:zeqn}
\end{aligned}\end{gathered}\end{equation}
where ${x}_{t}$ is period-t eXpenditure on the durable good and $\delta$ 
is the durable good's depreciation rate (a good with a lower value of 
$\delta $ is said to be ``more durable'').\opt{MarginNotes}{\marginpar{\tiny Note features of reality left out: no asymmetry between buying and selling the durable good (i.e. there is no restriction saying ${x}_{t}>0$.}}

The dynamic budget constraint is \opt{MarginNotes}{\marginpar{\tiny Only difference is must subtract expenditures on durables.}}
\begin{equation}\begin{gathered}\begin{aligned}
 {m}_{t+1} & =  ({m}_{t}-c_{t}-{x}_{t})\Rfree + y_{t+1}. \label{eq:xaccum}
\end{aligned}\end{gathered}\end{equation}

Bellman's equation is \opt{MarginNotes}{\marginpar{\tiny The reason $d_{t-1}$ is the state is that the level of durables in period $t$ is not determined until you choose spending on durables in that period.}}
\begin{equation}\begin{gathered}\begin{aligned}
        \vFunc_{t}({m}_{t},d_{t-1}) & =  \max_{\{c_{t},{x}_{t}\}} \left[\uFunc(c_{t},d_{t}) + \beta \vFunc_{t+1}({m}_{t+1},d_{t})\right],
\end{aligned}\end{gathered}\end{equation}
or, equivalently, \opt{MarginNotes}{\marginpar{\tiny Difference is that in second version $d_{t}$ is treated as control.}}
\begin{equation}\begin{gathered}\begin{aligned}
        \vFunc_{t}({m}_{t},d_{t-1}) & =  \max_{\{c_{t},d_{t}\}} \left[\uFunc(c_{t},d_{t}) + \beta \vFunc_{t+1}({m}_{t+1},d_{t})\right], \label{eq:bell}
\end{aligned}\end{gathered}\end{equation}
subject to
\begin{equation}\begin{gathered}\begin{aligned}
        {m}_{t+1} & =  \left({m}_{t}-c_{t}-\overbrace{(d_{t}-(1-\delta)d_{t-1})}^{={x}_{t}}\right)\Rfree+y_{t+1}
\end{aligned}\end{gathered}\end{equation}
or (substituting this into \eqref{eq:bell}),
\begin{equation}\begin{gathered}\begin{aligned}
        \vFunc_{t}({m}_{t},d_{t-1}) & =  \max_{\{c_{t},d_{t}\}} \left\{\uFunc(c_{t},d_{t}) + \beta \vFunc_{t+1}(({m}_{t}-c_{t}-(d_{t}-(1-\delta)d_{t-1}))\Rfree+y_{t+1},d_{t})\right\}. \nonumber
\end{aligned}\end{gathered}\end{equation}

Since this equation has two control variables, $c_{t}$ and $d_{t}$, there are two
first order conditions:

wrt $c_{t}$:
\begin{equation}\begin{gathered}\begin{aligned}
        \uFunc_{t}^{c} - \Rfree\beta \vFunc^{{m}}_{t+1}& =  0 \\
    \uFunc_{t}^{c}& =  \Rfree\beta \vFunc^{{m}}_{t+1}
\end{aligned}\end{gathered}\end{equation}

wrt $d_{t}$:
\begin{equation}
\uFunc_{t}^{d} = \beta(\Rfree \vFunc_{t+1}^{{m}} - \vFunc^{d}_{t+1}) = \Rfree\beta \label{eq:ud}
\vFunc_{t+1}^{{m}}-\beta \vFunc^{d}_{t+1}.
\end{equation}

Note that when taking the derivative with respect to $c_{t}$ you 
assume that $\partial d_{t}/\partial c_{t} = 0$ and vice versa.
Although the first order conditions will define a relationship
between the {\it optimal} values of $c_{t}$ and $d_{t}$, there is
no {\it mechanical} link that applies at this point.

Now we want to apply the Envelope theorem.  Basically, the Envelope 
theorem says that at the optimal levels of the control variables the 
partial derivative of the entire value function with respect to each 
control variable is zero.  This means that when taking the derivative 
with respect to a {\it state} variable you can simply ignore all terms 
that involve~$\partial c_{t}/ \partial {m}_{t}, \partial d_{t} / 
\partial {m}_{t}, \partial c_{t} / \partial d_{t-1}, $ and $\partial 
d_{t} / \partial d_{t-1}.$ So, for example, the full expression for 
the derivative of the value function with respect to ${m}_{t}$ is:
\begin{equation}\begin{gathered}\begin{aligned}
        \vFunc_{t}^{{m}} & =  \frac{\partial \uFunc(c_{t},d_{t})}{\partial c_{t}}
                        \frac{\partial c_{t}}{\partial {m}_{t}}
                  + \frac{\partial \uFunc(c_{t},d_{t})}{\partial d_{t}}
                        \frac{\partial d_{t}}{\partial {m}_{t}} \nonumber \\
              &  + \left[\frac{\partial {m}_{t+1}}{\partial {m}_{t}}
                  +       \frac{\partial {m}_{t+1}}{\partial c_{t}}
                              \frac{\partial c_{t}}{\partial {m}_{t}}
                  +       \frac{\partial {m}_{t+1}}{\partial d_{t}}
                              \frac{\partial d_{t}}{\partial {m}_{t}} 
                        \right] \beta \vFunc^{{m}}_{t+1} + \beta 
                              \vFunc^{d}_{t+1} \frac{\partial d_{t}}{\partial 
                              {m}_{t}} \label{eq:vmt}
\end{aligned}\end{gathered}\end{equation}
but the Envelope theorem tells us to ignore all the terms that involve 
$\partial c_{t}/\partial {m}_{t}$ or $\partial d_{t}/\partial {m}_{t}$; 
then because the only term in that whole mess above that does {\it 
not} involve either $\partial c_{t}/\partial {m}_{t}$ or $\partial 
d_{t}/\partial {m}_{t}$ is $\partial {m}_{t+1}/\partial {m}_{t} = \Rfree$ we 
have:
\begin{equation}\begin{gathered}\begin{aligned}
        \vFunc_{t}^{{m}} & =  \Rfree \beta \vFunc_{t+1}^{{m}}  \label{eq:vxtEqvxtp1}
\end{aligned}\end{gathered}\end{equation}

Applying the same Envelope theorem logic for $d_{t-1}$ yields:\footnote{Both here and in \eqref{eq:vmt}, the derivative $\vFunc^{d}_{t}$ should be understood as a differentiation with respect to $d_{t-1}$.}
\begin{equation}\begin{gathered}\begin{aligned}
        \vFunc_{t}^{d}  & =  \Rfree(1-\delta) \beta \vFunc^{{m}}_{t+1} \\
         & =  (1-\delta)  \Rfree \beta \vFunc_{t+1}^{{m}} \\
         & =  (1-\delta) \vFunc^{{m}}_{t} \label{eq:vdvsvx}
\end{aligned}\end{gathered}\end{equation}

Think now about the case where depreciation is 100 percent 
($\delta=1$); from \eqref{eq:vdvsvx} it is clear that in this case 
$\vFunc_{t}^{d} = 0$.  This makes sense because in this case the `durable' 
good is really a totally nondurable good.  $\vFunc^{d_{t-1}}_{t} = 0$ 
because the amount that you consumed of a nondurable good last period 
has no direct effect on your current utility (we have assumed that 
utility is time separable).

The $\delta= 0$ case is more interesting.  In this case the marginal 
utility of having an extra unit of durable good last period is equal 
to the marginal utility of having an extra unit of wealth this 
period.  Why?  Because if $\delta=0$ the durable good doesn't 
depreciate at all.  How much would it cost to buy another unit of 
durable good today?  One unit of wealth.  Because the durable does not 
depreciate from period to period and can be transformed into and out 
of wealth at a one-to-one price, it is exactly as valuable as a unit 
of wealth.

Now we want to try to derive a relationship between the contemporaneous
marginal utilities of $d$ and $c$.  From \eqref{eq:ud} we have:
\begin{equation}
\uFunc^{d}_{t} = \Rfree \beta \vFunc^{{m}}_{t+1} - \beta \vFunc^{d}_{t+1} \label{eq:beqns}.
\end{equation}
and \[ \Rfree \beta \vFunc^{{m}}_{t+1} = \uFunc^{c}_{t} \] and from 
\eqref{eq:vdvsvx} $\vFunc_{t+1}^{d} = (1-\delta) \vFunc^{{m}}_{t+1}$.  
Substituting these into (\ref{eq:beqns}):
\begin{equation}\begin{gathered}\begin{aligned}
         \uFunc_{t}^{d} & =   \uFunc_{t}^{c}-\beta(1-\delta) \vFunc_{t+1}^{{m}} \\
         & =  \uFunc_{t}^{c} - \frac{(1-\delta)}{\Rfree} \Rfree\beta \vFunc_{t+1}^{{m}}  \\
         & =  \left[1-\frac{(1-\delta)}{\Rfree}\right] \uFunc_{t}^{c} \\
         & =  \left[\frac{\rfree + \delta}{\Rfree}\right]   \uFunc_{t}^{c}\label{eq:uprimeeqn}
\end{aligned}\end{gathered}\end{equation}

Assuming $\delta<1$, this equation tells us that the marginal utility \opt{MarginNotes}{\marginpar{\tiny You're going to own this thing for a long time.  You don't buy a car because it is worth \$20,000 to you on the day you buy it; you buy a car because its {\it expected discounted} value to you over its lifetime is \$20,000 (or more).}}
{\it in the current period} of a unit of spending on the durable good 
is lower than the marginal utility of spending on the nondurable.  
Why?  Because the durable good will yield utility in the future as 
well as in the present.  What should be equated to the marginal 
utility of nondurables consumption is the total discounted lifetime 
utility from an extra unit of the durable good, not simply the 
marginal utility it yields right now.


Now assume the utility function is of the Cobb-Douglas form: $\uFunc(c,d)=\frac{({%
c^{1-\alpha }d^{\alpha })}^{1-\rho }}{1-\rho }.$  This implies that the
instantaneous marginal utilities with respect to $c$ and $d$ are:
\begin{equation}\begin{gathered}\begin{aligned}
        \uFunc^{c} & =  (c^{1-\alpha}d^{\alpha})^{-\rho} 
        (1-\alpha)c^{-\alpha}d^{\alpha}  \\
                   & =   (c^{1-\alpha}d^{\alpha})^{-\rho} (1-\alpha)(d/c)^{\alpha}
%\\                   & =   \left(c(d/c)^{\alpha}\right)^{-\rho} (1-\alpha)(d/c)^{\alpha}
%\\                   & =   c^{-\CRRA} (1-\alpha)(d/c)^{\alpha(1-\CRRA)}
%\\ \uFunc^{cc} & =   c^{-\CRRA} (1-\alpha)(\alpha(1-\CRRA))(d/c)^{\alpha(1-\CRRA)-1}(-dc^{-2})-\CRRA c^{-\CRRA-1}(1-\alpha)(d/c)^{\alpha(1-\CRRA)}
%\\ \uFunc^{cc} & =   c^{-\CRRA} (1-\alpha)(\alpha(1-\CRRA))(d/c)^{\alpha(1-\CRRA)-1}(-d/c)c^{-1})
%  \\ \uFunc^{cc} & =   c^{-\CRRA-1} (1-\alpha)(\alpha(1-\CRRA))(d/c)^{\alpha(1-\CRRA)-1}(-d/c)
%  \\ \uFunc^{cc} & =  -\CRRA(c^{0.5}d^{0.5})^{-\rho-1}0.5(d/c)^{0.5}+(c^{1-\alpha}d^{\alpha})^{-\rho} 0.5 (d/c)^{-0.5}(-d/c)/c
%  \\ \uFunc^{cd} & =  -\CRRA(c^{0.5}d^{0.5})^{-\rho-1}0.5(d/c)^{-0.5}(c^{-1}-d c^{-2})
\\        \uFunc^{d} & =  (c^{1-\alpha}d^{\alpha})^{-\rho} \alpha c^{1-\alpha} 
        d^{\alpha-1} \\
              & =  (c^{1-\alpha}d^{\alpha})^{-\rho}  \alpha (d/c)^{\alpha-1}
%\\ \uFunc^{dd} & =  -\CRRA(c^{0.5}d^{0.5})^{-\rho-1}0.5(c/d)^{0.5}+(c^{1-\alpha}d^{\alpha})^{-\rho} 0.5 (d/c)^{-0.5}c^{-1}
\end{aligned}\end{gathered}\end{equation}

Substituting these definitions into \eqref{eq:uprimeeqn} gives:
\begin{equation}\begin{gathered}\begin{aligned}
(c^{1-\alpha}d^{\alpha})^{-\rho} \alpha (d/c)^{\alpha-1}& = 
(c^{1-\alpha}d^{\alpha})^{-\rho}  (1-\alpha) 
(d/c)^{\alpha}\left(\frac{\rfree+\delta}{\Rfree}\right) \\
        \frac{\alpha}{1-\alpha} & =  (d/c)\left(\frac{\rfree+\delta}{\Rfree}\right)  \\
        d/c & =  \left(\frac{\alpha}{1-\alpha}\right)\left( \frac{\Rfree}{\rfree+\delta} \right)\equiv \gamma \label{eq:gamma}
\end{aligned}\end{gathered}\end{equation}

What this implies is that whenever the level of nondurables 
consumption changes, the level of the {\it stock} of durables should 
change by the same proportion.  Because {\it expenditures} on durable 
goods are equal to the {\it change} in the stock plus depreciation, a 
change in $c$ implies spending on durables large enough to immediately 
adjust the stock to the new target level.  (Recall that $d_{t}$ was 
the {\it stock} of durable good owned in period $t$, while {\it 
spending} on the durable good was defined as ${x}_{t} = d_{t} - 
(1-\delta) d_{t-1}$.)

Define $\gamma = d_{t}/c_{t}$ as in \eqref{eq:gamma}.  Now consider a
consumer who had consumed the same amount of the nondurable good for
periods $c_{t-2} = c_{t-1}$ but who between period $t-1$ and period
$t$ learns some good news about permanent income; she adjusts her
nondurables consumption up so that $c_{t}/c_{t-1} = (1+\epsilon_{t})$.
This implies that the level of {\it spending} in period $t$ is:
\begin{equation}\begin{gathered}\begin{aligned}
        {x}_{t}   & =  d_{t} - (1-\delta) d_{t-1} = \gamma [c_{t} - (1-\delta) 
        c_{t-1}]\\
        {x}_{t-1} & =  % d_{t-1} - (1-\delta) d_{t-2} \\                 & =  
\gamma [c_{t-1} - (1-\delta) c_{t-2}]  \\
                & =  \gamma \delta c_{t-1}  \\
        {x}_{t}/{x}_{t-1} & =  \gamma[c_{t-1}(1+\epsilon_{t}) - 
        (1-\delta)c_{t-1}]/ \gamma \delta c_{t-1}  \\
         & =  \frac{\epsilon_{t} + \delta}{\delta} \label{eq:dursspending}
\end{aligned}\end{gathered}\end{equation}

Assuming $\delta < 1$, this equation implies that {\it spending on \opt{MarginNotes}{\marginpar{\tiny Briefly discuss meaning: your house is 3 times your perm income; so fluctuations in perm income mean you will be perpetually buying new chunks of house (when P goes up) or selling off chunks of your house (when P) goes down.  Obviously missing transactions costs; take class next year to learn about those!  (Also point out that durables spending is in fact much more volatile than nondurables.)}}
durable goods should be more variable than spending on nondurable 
goods.} For goods with a low depreciation rate, spending should be 
much more variable.  This is true because the ratio of the stock of 
durables to income is much larger than the ratio of the average level 
of spending on durables to income.

A further implication of this model is that the degree of correlation between nondurables spending growth and durables spending growth
depends on the frequency under consideration.  For a given quarterly depreciation rate (say, 5 percent per quarter), the durable good will have almost completely depreciated over the course of 10 years $=$ 40 quarters because $0.95^{40}=0.12$.  According to the model, over an interval long enough for the durable to have completely depreciated, the rate of growth of spending on the durable should match the rate of growth of spending of the nondurable, because over such a long interval they are really both nondurable.

Some evidence on this proposition is provided in the Jupyter notebook available at \href{https://github.com/llorracc/Jupyter/blob/master/notebooks/Durables-vs-Nondurables-At-Low-And-High-Frequencies.ipynb}{here}.
% or runnable by \href{https://mybinder.org/v2/gh/llorracc/Jupyter/master?filepath=notebooks\%2FDurables-vs-Nondurables-10y-vs-1q.ipynb}{mybinder}.
% As of 20181206, mybinder disabled because requires too many things (matplotlib, seaborne, pandas) that are not part of requirements.txt and are BIG

 
\input handoutBibMake


\end{document}
